\documentclass[12pt,a4paper]{article}
\usepackage{fontsetup}
\usepackage[english, russian, ukrainian]{babel}


\usepackage{geometry}
\geometry{left=3cm,right=1cm,top=1cm,bottom=1cm}
\usepackage{xcolor}
\usepackage{graphicx}
\usepackage{tabularray}
\UseTblrLibrary{booktabs}
\usepackage[
    colorlinks=true,
	urlcolor = blue, %Colour for external hyperlinks
	linkcolor  = malina, %Colour of internal links
	citecolor  = green, %Colour of citations
	bookmarks = true,
	bookmarksnumbered=true,
	unicode,
	linktoc = all,
	hypertexnames=false,
	pdftoolbar=false,
	pdfpagelayout=TwoPageRight,
	pdfauthor={Ponomarenko S.M. aka sergiokapone},
	pdfdisplaydoctitle=true,
	pdfencoding=auto
	]%
	{hyperref}
		\makeatletter
	\AtBeginDocument{
	\hypersetup{
		pdfinfo={
		Title={\@title},
		}
	}
	}
	\makeatother



\begin{document}

\begin{center}
    \textbf{\large Висновок відповідального редактора}\\[0.5em]
    \textbf{щодо рукопису «Електрика та магнетизм»}\\[0.2em]
    \textbf{автора Пономаренка С.М.}
\end{center}

\vspace{1.5em}

\noindent Цитати, посилання та фактичний матеріал перевірено.

\vspace{0.5em}
\noindent Внаслідок перевірки рукопису на текстові збіги встановлено 9,2\% збігів: кількість ідентичних слів з домашньої бази складає 2,69\%, кількість ідентичних слів з інтернету -- 4,81\%, а також кількість ідентичних слів з програм обміну базами даних -- 1,7\%, що пояснюється використанням слів та словосполучень термінологічного характеру.

\vspace{0.5em}
\noindent У результаті перевірки рукопису Електрика та магнетизм» автора Пономаренка С. М. встановлено, що:
\begin{itemize}
    \item[--] рукопис відповідає силабусу дисципліни;
    \item[--] якість текстового та графічного матеріалу в рукописі на високому рівні;
    \item[--] написання науково-технічних термінів, одиниць вимірювання, а також наведені у рукописі символи та позначення відповідають нормативним документам;
    \item[--] рукопис оформлено із дотриманням вимог Порядку підготовки та вимоги до рукописів навчальних видань в НН ФТІ для надання грифів КПІ ім. Ігоря Сікорського, затвердженого Вченою радою НН ФТІ (протокол №X/X від XX.XX.2026 р.).
\end{itemize}

\vspace{4em}

\noindent
\textbf{Відповідальний редактор}:\\
перший заступник директора НН ФТІ\\ КПІ ім. Ігоря Сікорського,\\
доцент, к.т.н. \hfill Тетяна ЛИТВИНОВА



\vspace{3em}

\noindent Розглянуто на засіданні кафедри прикладної фізики\\
(протокол №10-2026 від «21» жовтня 2026 р.)

\end{document}